% Pursuing Stacks (cotes 134-2 to 134-8): archive cote 134-7, site page 95 (PDF page 96), typescript page 537.
% typescript. Draft compared: ps6.tex:4772-4843.
% AI-transcribed by Claude Opus 5.5 (claude-opus-5-5, adaptive thinking, effort high), Message Batches API, batch ps4 (msgbatch_012Aoc1NuskATPsyBeKSYK7J), 2026-11-12; not yet checked by eye.
\scanpage{134-7/95}
\mspage{537}
is a weak equivalence too. \hand{Thus a ``schematic homotopy type over $k$''
should define a functor $Alg_{/k} \to Hot$.}
For any choice of a``section''
\[ (3) \qquad s \in X_0(k) \]
of $X_*$ , Illusie's constructions yield sheaves
\[ (4) \qquad \underline{\Pi}_i(X_*,s) \qquad ( \; i \geq 0 \; )\hand{,} \]
where $\underline{\Pi}_0$ is a sheaf of sets, $\underline{\Pi}_1$ a sheaf of groups, \inserted[typed]{acting on} \xout{and} the higher
$\underline{\Pi}_i$'s \inserted[typed]{which\hand{,} \xout{are}} are abelian sheaves. We will be mainly interested of course in the
case when $\underline{\Pi}_0$ is the final sheaf (we'll say that $X_*$ is \uline{relatively O-}
\uline{connected over} k), and moreover $\underline{\Pi}_1$ is abelian and it's action on
the higher $\underline{\Pi}_i$ is trivial (let's say in this case that the relative
homotopy type defined by the semisimplicial sheaf $X_*$ is \uline{pseudo}-\uline{abelian}).
In this case, up to canonical isomorphisms, the \xout{\unsure{sequence of}} abelian
sheaves $\underline{\Pi}_i(X_*,s)$ do not depend on the choice of s, and (provided the
pseudoabelian notion is defined locally) they make a sense, independent%
-ly \inserted{even} of the existence \deleted{even} of a section s. Our hope is, by suitably
restricting our notion of ``bundle'', to get on the abelian sheaves $\underline{\Pi}_i$
for $i \geq 2$ , or even on all $\underline{\Pi}_i$ $(i \geq 1)$, a natural structure of an
$\underline{O}_k$-module. \xout{maybe} and one moreover which in many ``good'' cases turns it
into a quasi-coherent $\underline{O}_k$-module. Also, the natural maps
\[ (5) \qquad \hand{$\pi$}_i(X_*(k)) \to \underline{\Pi}_i(X_*)(k) \]
should be isomorphisms\hand{,} i.e. taking $\hand{$\pi$}_i$ and $\underline{\Pi}_i$ should commute to taking
sections (which will imply indeed that if (1) is a quasi-isomorphism
then so is (2)). Thus a first basic question here is to find a suitable
notion of a ``bundle'' over k, in such a way that for semisimplicial
bundles satisfying some mild extra assumption (such as $X_0=X_1=e$) implying
that $X_*$ is pseudoabelian, the maps (5) should be isomorphisms. I ex-%
pect this to be true for the notion we have been working with so far,
namely for ``unipotent bundles'' (cf. section 111), but I hav'nt made
any attempt yet to prove this. But even granting this property for
a given notion of ``bundles'', the Module structure on the $\underline{\Pi}_i$'s at this
point remains still a mystery,\unsure{\hand{$^*$}} as long as we do'nt tie them in with the
Lie functor (compare section 118) ...

\uline{Remark}. Maybe, in the set-up of a general topos \inserted{T} and taking section\unsure{\hand{s}} on
the latter, the maps (5) are isomorphisms, whenever the \deleted{cohomology}
\deleted{g}

\margin{95}
